% !TEX program = xelatex
% Статья для сборника НИТУ «МИСИС» (курс «Научное письмо»).
%
% Правила текста:
%   * ссылка на источник — \cite с ключом из paper/refs.bib основного проекта;
%   * утверждение из источника — комментарий % C### в конце строки (id из paper/claims.csv);
%   * список литературы не набирается руками: python scripts/bib_to_tex.py <путь>/article.tex;
%   * числа о наших экспериментах — только из experiments/results (в этой версии их нет).

\documentclass{misisarticle}

\usepackage{algorithm}
\usepackage[noend]{algpseudocode}
\floatname{algorithm}{Алгоритм}
\algrenewcommand\algorithmicrequire{\textbf{Вход:}}
\algrenewcommand\algorithmicensure{\textbf{Выход:}}
\algrenewcommand\algorithmicif{\textbf{если}}
\algrenewcommand\algorithmicthen{\textbf{то}}
\algrenewcommand\algorithmicelse{\textbf{иначе}}
\algrenewcommand\algorithmicloop{\textbf{цикл}}
\algrenewcommand\algorithmicreturn{\textbf{вернуть}}
\algrenewcommand\algorithmicfor{\textbf{для}}
\algrenewcommand\algorithmicdo{\textbf{выполнять}}

% --- Классификационные коды (TODO: сверить с рубрикаторами) ---
\grnti{28.23}
\asjc{1702, 1710}
\oecd{1020}

% --- Название ---
\title{Агентный RAG-конвейер с планированием запроса и контролем достаточности контекста для вопросно-ответного поиска по технической документации}

% --- Авторы: \misisauthor{ФИО}{аффилиация}{город, страна}{email} ---
% Дерзновенко Александр Александрович. TODO: уточнить кафедру / институт.
\misisauthor{А.~А.~Дерзновенко}{НИТУ «МИСИС»}{Москва, Россия}{m2306033@edu.misis.ru}

% --- Аннотация и ключевые слова ---
\abstract{%
  Генерация, дополненная поиском (RAG), в базовом варианте выполняет поиск один раз по исходной формулировке вопроса и не проверяет, хватает ли найденного для ответа. При недостаточном контексте языковые модели склонны скорее отвечать, чем воздерживаться, и система отвечает по неполному или нерелевантному контексту вместо повторного поиска или отказа. В статье предложен агентный RAG-конвейер для вопросно-ответного поиска по технической документации, не требующий обучения моделей. Модуль оценки достаточности контекста проверяет покрытие каждой части вопроса, агент-планировщик по списку непокрытых частей выбирает стратегию следующего раунда поиска (прямой поиск, переформулирование или декомпозиция) и формирует запросы только для этих частей, а память сессии исключает уже просмотренные фрагменты. При исчерпании бюджета итераций система отказывается от ответа и указывает, какие части вопроса не покрыты; сгенерированный ответ проверяет отдельный агент-критик. Проведён анализ близких методов управления поиском, дана формальная постановка задачи с критериями качества ответов, отказов и стоимости, описаны компоненты конвейера, правило завершения и реализация на LangGraph. Экспериментальная проверка в статью не входит: определены её критерии и схема.
}

\keywords{%
  генерация, дополненная поиском, агентный RAG, оценка достаточности контекста, обоснованный отказ, агент-критик, техническая документация, большие языковые модели
}

\begin{document}

\maketitle

% ============================================================
\section{Введение}

Генерация, дополненная поиском (retrieval-augmented generation, RAG), дополняет большую языковую модель (LLM) поиском по внешней базе знаний; в простейшем варианте конвейер состоит из индексации, поиска и генерации и называется схемой «Retrieve-Read»~\cite{gao2023retrieval}. % C116
В этой работе RAG рассматривается как основа ассистента технической поддержки, который отвечает на вопросы пользователей по документации продукта.
Такие данные есть и в открытых бенчмарках: набор TechQA в RAGBench составлен из запросов пользователей с технических форумов IBM~\cite{friel2024ragbench}, % C110
а набор EManual содержит вопросы о телевизоре Samsung Smart TV и его пульте~\cite{friel2024ragbench}. % C104

Базовый RAG извлекает фиксированное число документов независимо от того, нужен ли поиск~\cite{asai2023selfrag}. % C003
На этапе поиска он выбирает нерелевантные фрагменты и может терять важную информацию, а для сложного вопроса одного поиска по исходному запросу может не хватить~\cite{gao2023retrieval}. % C117, C119
На этапе генерации модель может выдать содержание, не подкреплённое найденным контекстом~\cite{gao2023retrieval}. % C118
При недостаточном контексте модели в целом скорее галлюцинируют, чем воздерживаются от ответа~\cite{joren2024sufficient}. % C147
На бенчмарке CRAG ни одно из рассмотренных авторами RAG-решений не превысило 20~\% по показателю truthfulness, поскольку все они вносят галлюцинации из нерелевантных результатов поиска~\cite{yang2024crag}. % C249
Лучшие (SOTA) решения достигают более высокого truthfulness, но доля галлюцинаций у них остаётся от 16 до 25~\%~\cite{yang2024crag}. % C257

Исходное предположение работы состоит в том, что на технической документации эти недостатки проявляются сильнее.
Вопросы пользователей бывают длинными и содержат версии продукта, симптомы и фрагменты логов, поэтому поиск по исходной формулировке может промахиваться мимо нужного фрагмента, а в руководствах пользователя нужная инструкция может теряться среди похожих.
Проверка этого предположения относится к экспериментальной части работы и в статью не входит.

Известные улучшения переносят решение о поиске внутрь конвейера: адаптивный поиск позволяет системе самой определять, нужен ли внешний поиск и когда его остановить~\cite{gao2023retrieval}. % C125
Во многих таких методах решение принимают обучаемые компоненты: Self-RAG обучает языковую модель поиску по требованию~\cite{asai2023selfrag}, % C005
в Corrective RAG дообучается внешний оценщик релевантности~\cite{yan2024corrective}, % C064
в Adaptive-RAG --- классификатор сложности запроса~\cite{jeong2024adaptive}. % C028
Агентность при этом имеет цену: в экспериментах Ferrazzi et al. агентная конфигурация RAG оказалась дороже до 3,6 раза из-за дополнительных шагов рассуждения и повторных вызовов инструмента, а ни улучшенный конвейер (Enhanced RAG), ни агентный не оказался универсально лучшим~\cite{ferrazzi2026agentic}. % C092, C093

Сигнал для решения «искать дальше или отвечать» дают работы о достаточности контекста.
Контекст называют достаточным, если он содержит всю информацию, необходимую для определённого ответа на вопрос~\cite{joren2024sufficient}; % C142
в разделе об ограничениях авторы этого определения указывают, что могли бы использовать свой автооценщик, чтобы итеративно решать, искать ещё или отвечать~\cite{joren2024sufficient}. % C151
Итеративные методы FAIR-RAG~\cite{asl2025fairrag}, S2G-RAG~\cite{li2026s2grag}, SIM-RAG~\cite{yang2025knowing} и PAR$^2$-RAG~\cite{li2026par2} реализуют подобный цикл, но различаются тем, обучаются ли их компоненты и что система делает, когда свидетельств так и не хватило (раздел~\ref{sec:related}).

Цель работы --- разработать агентный RAG-конвейер, в котором дополнительные шаги поиска запускаются только при недостаточном контексте, а при исчерпании бюджета итераций система отказывается от ответа.
Для этого в статье решаются следующие задачи:
\begin{enumerate}[nosep]
  \item спроектировать архитектуру конвейера из четырёх компонентов: агента-планировщика, модуля оценки достаточности контекста по частям вопроса, агента-критика и памяти сессии;
  \item формализовать правила повторного поиска, остановки по бюджету итераций и отказа от ответа, а также входы и выходы планировщика, выбирающего стратегию поиска;
  \item реализовать конвейер, отделив модель-судью от модели-генератора;
  \item определить критерии и схему проверки работоспособности конвейера относительно базового RAG по качеству ответов, корректности отказов и стоимости.
\end{enumerate}
Сама экспериментальная проверка --- предмет следующего этапа работы; в статье определены её критерии и схема.

Вклад работы состоит в следующем.
\begin{enumerate}[nosep]
  \item Цикл управления поиском без обучения моделей, в котором оценка достаточности контекста по частям вопроса определяет не только, нужно ли искать ещё, но и как: по списку непокрытых частей планировщик выбирает стратегию следующего раунда (прямой поиск, переформулирование или декомпозиция) и формирует запросы только для этих частей, а память сессии исключает уже просмотренные фрагменты.
  \item Обоснованный отказ как явный исход цикла, для которого определены метрики: при исчерпании бюджета итераций система отказывается от ответа и указывает непокрытые части вопроса, а сгенерированный ответ дополнительно проверяет отдельный агент-критик. Оба механизма реализованы промптами, без обучаемых оценщиков и критиков.
  \item Постановка задачи и схема проверки, в которых качество ответов и отказов оценивается вместе с ценой агентности --- задержкой, числом токенов и раундов поиска --- относительно базового RAG с теми же параметрами поиска.
\end{enumerate}

Дальнейшее изложение построено так.
В разделе~\ref{sec:related} проанализированы методы управления поиском в RAG и подходы к их оценке.
В разделе~\ref{sec:problem} дана постановка задачи.
В разделе~\ref{sec:method} описаны компоненты конвейера, правило завершения и реализация.

% ============================================================
\section{Анализ существующих работ}\label{sec:related}

\subsection{Парадигмы RAG и преобразование запроса}

Gao et al. различают наивный, продвинутый и модульный RAG~\cite{gao2023retrieval}.
В продвинутом RAG на стадии до поиска основное внимание уделяется структуре индекса и исходному запросу~\cite{gao2023retrieval}; % C120
к типовым методам этой стадии относят переписывание, преобразование и расширение запроса~\cite{gao2023retrieval}. % C121
Модульный RAG не ограничен последовательной схемой «поиск --- генерация» и включает итеративный и адаптивный поиск~\cite{gao2023retrieval}. % C122
Итеративный поиск многократно обращается к базе знаний на основе исходного запроса и уже сгенерированного текста~\cite{gao2023retrieval}, % C124
но может страдать от семантической разорванности и накопления нерелевантной информации~\cite{gao2023retrieval}. % C126

Ma et al. указывают на неизбежный разрыв между входным текстом и знанием, которое нужно найти~\cite{ma2023query}. % C179
В простейшем варианте их подхода запрос переписывает LLM по few-shot-промпту, и результатом бывает ни одного, один или несколько запросов~\cite{ma2023query}. % C180
На HotpotQA поиск по исходному сложному вопросу ухудшал результат модели-читателя, что авторы связывают с шумом в найденном~\cite{ma2023query}. % C185
Вместе с тем Wang et al. исходят из того, что не каждому запросу нужен поиск~\cite{wang2024searching}, % C217
и сообщают, что наивысшие оценки поиска дало сочетание HyDE и гибридного поиска, а переписывание и декомпозиция запроса не улучшили поиск так же эффективно~\cite{wang2024searching}. % C256
Выводы этих работ о пользе преобразования запроса различаются, поэтому в предлагаемом конвейере стратегия не фиксируется заранее, а выбирается на каждом раунде по результату проверки достаточности.

Singh et al. описывают агентный RAG как смену парадигмы, при которой в конвейер вводятся автономные агенты, способные принимать решения динамически~\cite{singh2025agentic}. % C206
Управляющий слой такого конвейера решает, когда извлекать, как переформулировать запрос и когда собрано достаточно свидетельств~\cite{singh2025agentic}. % C207
Среди паттернов агентов выделяются планирование --- самостоятельное разбиение сложной задачи на подзадачи --- и рефлексия --- итеративная оценка и улучшение собственных выходных данных~\cite{singh2025agentic}. % C208, C209
Авторы обзора отмечают, что агентность вносит сложность координации, задержку и вычислительные затраты~\cite{singh2025agentic}, % C214
а эффективные системы описывают как сочетающие автономию с ограниченным горизонтом планирования, заданными политиками доступа к инструментам и явными критериями остановки~\cite{singh2025agentic}. % C216

\subsection{Управление поиском с обучаемыми компонентами}

Self-RAG по входу и предыдущей генерации определяет, полезен ли поиск, и при необходимости выдаёт токен поиска~\cite{asai2023selfrag}. % C001
Дообученная модель использует critique-токены, чтобы оценивать собственные предсказания после каждого сгенерированного сегмента~\cite{asai2023selfrag}; % C002
разметка этих токенов для обучения критика получена от GPT-4 и дистиллирована в собственную модель-критика~\cite{asai2023selfrag}. % C011
Релевантность фрагмента в Self-RAG оценивается относительно всего входа~\cite{asai2023selfrag}. % C008
Оценка по отдельным частям вопроса в работе не описана.

FLARE не требует дополнительного обучения и применим к существующим языковым моделям на этапе вывода~\cite{jiang2023active}. % C129
Повторный поиск запускается, если вероятность любого токена временного следующего предложения ниже порога~\cite{jiang2023active}. % C131
Авторы признают, что наивное чередование генерации и поиска увеличивает накладные расходы и стоимость генерации~\cite{jiang2023active}. % C140

Corrective RAG оценивает релевантность каждого документа запросу по отдельности: пара «вопрос --- документ» получает одну оценку~\cite{yan2024corrective}. % C052
Оценщик --- модель T5-large~\cite{yan2024corrective}. % C055
При оценке выше верхнего порога выбирается действие Correct, ниже нижнего --- Incorrect, иначе --- Ambiguous~\cite{yan2024corrective}; % C056
при Incorrect система обращается к новым источникам знаний~\cite{yan2024corrective}. % C057
Авторы отмечают, что эффективность метода сильно зависит от точности оценщика релевантности~\cite{yan2024corrective}, % C058
а его дообучение считают неизбежным~\cite{yan2024corrective}. % C064

Adaptive-RAG определяет сложность запроса до попытки ответа, и стратегия выбирается до выполнения~\cite{jeong2024adaptive}. % C027
Классификатор сложности --- меньшая языковая модель, обученная относить запрос к одному из трёх уровней~\cite{jeong2024adaptive}. % C028
Авторы признают, что автоматическая разметка для его обучения может ошибочно назначать метки сложности~\cite{jeong2024adaptive}. % C043

В этих методах решение о поиске опирается на свойство запроса, документа или генерируемого текста.
Проверки того, какие именно части вопроса покрыты найденным контекстом, в их описаниях не найдено.

\subsection{Итеративный поиск по сигналу достаточности}

Joren et al. определяют метку достаточности по паре «вопрос --- контекст», без эталонного ответа~\cite{joren2024sufficient}, % C143
и сообщают о точности автооценщика достаточности 93~\%~\cite{joren2024sufficient}; % C144
эта оценка получена на размеченном наборе из 115 примеров, лучший результат дал Gemini 1.5 Pro в режиме 1-shot~\cite{joren2024sufficient}. % C253
Авторы отмечают, что простое воздержание при недостаточном контексте --- грубый подход, который может снизить общее качество~\cite{joren2024sufficient}. % C149
Поэтому в предлагаемом конвейере отказ --- не реакция на первую же оценку недостаточности, а исход после исчерпания бюджета повторного поиска.
Цена такого решения --- отказы на вопросах, на которые модель ответила бы верно, --- измеряется отдельно (раздел~\ref{sec:problem}).

FAIR-RAG считает доказательства достаточными только при подтверждении всех требуемых находок~\cite{asl2025fairrag}. % C068
При недостаточности агент уточнения учитывает уже подтверждённые факты, чтобы новые запросы были точнее и не повторяли прежние поиски~\cite{asl2025fairrag}; % C070
цикл ограничен тремя итерациями ради баланса между полнотой и задержкой~\cite{asl2025fairrag}. % C069
Модуль оценки достаточности SEA, в отличие от методов на основе дообучения, не требует обучения модели~\cite{asl2025fairrag}. % C072
О нехватке доказательств система сообщает по инструкции в промпте финальной генерации~\cite{asl2025fairrag}. % C071
Авторы признают, что каждый цикл уточнения увеличивает задержку и вычислительные затраты~\cite{asl2025fairrag}, % C075
а среди ошибок выделяют преждевременное суждение о достаточности~\cite{asl2025fairrag}. % C076

S2G-RAG на каждом шаге предсказывает, достаточно ли накопленных свидетельств, и при недостаточности выдаёт структурированные описания недостающей информации~\cite{li2026s2grag}. % C166
Ответчик вызывается либо когда принято решение отвечать, либо когда достигнут максимальный бюджет~\cite{li2026s2grag}. % C168
Судья дообучается на метках, которые размечает более сильная модель-учитель~\cite{li2026s2grag}. % C173

В PAR$^2$-RAG планировщик порождает набор подзапросов, покрывающих разные аспекты информационной потребности~\cite{li2026par2}, % C154
а решение о продолжении поиска или остановке принимает отдельный контроллер достаточности~\cite{li2026par2}. % C155
Если контроллер выбирает остановку, цикл завершается и последний промежуточный ответ возвращается как итоговый~\cite{li2026par2}. % C156
Решение контроллера принимает LLM, управляемая промптом, и так же через промпт генерируется ответ~\cite{li2026par2}. % C255, C254
Авторы отмечают, что методы с планированием могут давать широкий, но статичный набор запросов, который перестаёт соответствовать меняющимся промежуточным данным~\cite{li2026par2}. % C158

SIM-RAG формулирует задачу многораундового RAG как умение продолжать поиск только при необходимости или воздерживаться от ответа~\cite{yang2025knowing}. % C230
Шаги «ответ --- проверка --- поиск» повторяются, пока критик не признает ответ достаточным~\cite{yang2025knowing}, % C232
и система спроектирована так, чтобы воздерживаться от ответа при недостаточности свидетельств~\cite{yang2025knowing}. % C234
Параметры базовой LLM не меняются~\cite{yang2025knowing}, % C233
но критик обучается на траекториях поиска, размеченных как успешные, если они приводят к правильному ответу~\cite{yang2025knowing}. % C231
Основные вычислительные затраты приходятся на генерацию данных~\cite{yang2025knowing}, % C236
а для сохранения качества критику может потребоваться переобучение~\cite{yang2025knowing}. % C237

\subsection{Цена агентности и оценка RAG}

По наблюдению Ferrazzi et al., основной вклад в задержку агентного RAG вносят вызовы LLM~\cite{ferrazzi2026agentic}. % C094
Агент в их экспериментах выполнял повторный поиск в 10~\% случаев, и 53~\% найденных при этом документов оставались теми же~\cite{ferrazzi2026agentic}. % C098
Значит, повторный поиск без учёта уже просмотренного больше чем наполовину повторяет прежний результат; память сессии в предлагаемом конвейере исключает такие повторы.

Отказ от ответа выделяется как отдельная способность RAG-систем: Gao et al. называют её Negative Rejection --- умение воздержаться, когда найденные документы не содержат нужных знаний~\cite{gao2023retrieval}. % C127
Те же авторы отмечают, что метрики оценки RAG пока не образуют зрелого стандартизированного подхода~\cite{gao2023retrieval}. % C128
Ragas определяет верность ответа контексту через возможность вывести утверждения ответа из найденного контекста~\cite{es2024ragas} % C079
для этого LLM сначала разбивает ответ на более короткие утверждения~\cite{es2024ragas}. % C080
ARES требует человеческой валидационной выборки примерно из 150 размеченных примеров~\cite{saadfalcon2024aresautomated}. % C195

Бенчмарк CRAG (не путать с методом Corrective RAG) содержит 4\,409 пар «вопрос --- ответ»~\cite{yang2024crag}. % C241
В задаче Task 3 кандидатами для поиска служат 50 веб-страниц на вопрос~\cite{yang2024crag}. % C245
При автоматической оценке неверные ответы штрафуются, верные поощряются, а пропуск («не знаю») получает ноль~\cite{yang2024crag}; % C247
оценки двух LLM-судей согласуются с ручной разметкой со средним F1 94,7 и 98,9~\%~\cite{yang2024crag}. % C248
При такой шкале отказ выгоднее неверного ответа, но она не отличает обоснованный отказ от необоснованного, поэтому в разделе~\ref{sec:problem} шкала дополняется метриками отказа.

\subsection{Итоги анализа}

Сравнение методов приведено в табл.~\ref{tab:related}.

\begin{table*}[t]
  \caption{Сравнение методов управления поиском}
  \label{tab:related}
  \centering
  \footnotesize
  \begin{tabular}{@{}>{\raggedright\arraybackslash}p{2.2cm}>{\raggedright\arraybackslash}p{3.4cm}>{\raggedright\arraybackslash}p{4.3cm}>{\raggedright\arraybackslash}p{2.6cm}>{\raggedright\arraybackslash}p{4.3cm}@{}}
    \toprule
    \textbf{Метод} & \textbf{Обучаемые компоненты} & \textbf{Что запускает (повторный) поиск} & \textbf{Проверка ответа} & \textbf{Исход при нехватке свидетельств} \\
    \midrule
    Self-RAG~\cite{asai2023selfrag} & генератор с reflection-токенами & токен поиска по входу и предыдущей генерации & critique-токены внутри генерации & не описан$^{*}$ \\ % C005, C001, C002
    Corrective RAG~\cite{yan2024corrective} & оценщик релевантности T5-large & оценка пары «вопрос --- документ» ниже нижнего порога & нет$^{*}$ & обращение к новым источникам знаний \\ % C064, C055, C052, C056, C057
    Adaptive-RAG~\cite{jeong2024adaptive} & классификатор сложности & уровень сложности, определённый до ответа & нет$^{*}$ & не описан$^{*}$ \\ % C028, C027
    FAIR-RAG~\cite{asl2025fairrag} & нет (оценщик достаточности SEA не обучается) & не подтверждены все требуемые находки; не более трёх итераций & нет$^{*}$ (достаточность проверяется до генерации) & инструкция в промпте генерации сообщить о нехватке доказательств \\ % C072, C068, C069, C071
    S2G-RAG~\cite{li2026s2grag} & судья достаточности & метка недостаточности и структурированные пробелы & нет$^{*}$ & по достижении бюджета вызывается ответчик; отказ не описан$^{*}$ \\ % C173, C166, C168
    PAR$^2$-RAG~\cite{li2026par2} & нет: контроллер и генератор ответа --- LLM с промптами & план подзапросов и контроллер достаточности & нет$^{*}$ & при остановке возвращается последний промежуточный ответ; отказ не описан$^{*}$ \\ % C255, C254, C154, C155, C156
    SIM-RAG~\cite{yang2025knowing} & внешний критик (LLM не меняется) & критик не принял ответ & критик & отказ от ответа \\ % C231, C233, C232, C234
    \midrule
    Предлагаемый & нет & непокрытые части вопроса; стратегию выбирает планировщик; не более $B$ раундов & агент-критик & обоснованный отказ с перечнем непокрытых частей \\
    \bottomrule
  \end{tabular}

  \vspace{2pt}
  \raggedright\footnotesize $^{*}$~В работе механизм не описан (вывод по отсутствию упоминания, а не утверждение авторов).
\end{table*}

Сравнение указывает на два отличия предлагаемого метода от рассмотренных.
Первое касается отказа.
Как отдельный исход цикла он заявлен только в SIM-RAG, и там решение принимает обучаемый критик~\cite{yang2025knowing}. % C234, C232, C231
FAIR-RAG сообщает о нехватке доказательств по инструкции в промпте финальной генерации~\cite{asl2025fairrag}; % C071
в S2G-RAG по достижении бюджета вызывается ответчик~\cite{li2026s2grag}, % C168
а в PAR$^2$-RAG при остановке возвращается последний промежуточный ответ~\cite{li2026par2}. % C156
Отдельного исхода-отказа в описаниях этих трёх методов не найдено.

Второе отличие касается способа повторного поиска.
FAIR-RAG и S2G-RAG уже используют сведения о недостающем: в FAIR-RAG новые запросы учитывают подтверждённые факты~\cite{asl2025fairrag}, % C070
а судья S2G-RAG выдаёт структурированные описания недостающей информации~\cite{li2026s2grag}. % C166
Однако выбора между прямым поиском, переформулированием и декомпозицией в их описаниях не найдено; в предлагаемом методе такой выбор делает планировщик по списку непокрытых частей вопроса.

Наконец, в рассмотренных работах близкие методы оцениваются на вопросах общего домена: FAIR-RAG --- в том числе на HotpotQA~\cite{asl2025fairrag}, % C073
S2G-RAG использует в качестве корпуса дампы Wikimedia~\cite{li2026s2grag}, % C174
SIM-RAG --- TriviaQA, HotpotQA и 2Wiki~\cite{yang2025knowing}. % C240
Оценок этих методов на технической документации в рассмотренных работах не найдено.

% ============================================================
\section{Постановка задачи}\label{sec:problem}

\subsection{Обозначения}

Пусть $\mathcal{D}$ --- корпус фрагментов документации, $q$ --- вопрос пользователя, $\tau$ --- время, на которое задан вопрос.
Модуль поиска $R(s, k, X)$ возвращает $k$ наиболее близких к запросу $s$ фрагментов корпуса, не входящих в множество $X \subseteq \mathcal{D}$.
Генератор $G$ --- языковая модель с фиксированными параметрами.
По вопросу $q$ система возвращает исход $y$: либо ответ $a$, либо отказ $\bot$; отказ по бюджету сопровождается обоснованием $\rho$.

Вопрос разбивается на части --- минимальные информационные потребности:
\begin{equation}
  P(q) = \{p_1, \dots, p_m\}, \quad 1 \le m \le 4,
  \label{eq:parts}
\end{equation}
причём простой вопрос состоит из одной части.
После $t$ раундов поиска накоплен контекст $E_t \subseteq \mathcal{D}$, $|E_t| \le N_{\max}$.
Покрытие части $c_t(p_i) = 1$, если $E_t$ явно содержит сведения, необходимые для ответа на $p_i$ (для частей, зависящих от времени, --- сведения, верные на момент $\tau$), иначе $c_t(p_i) = 0$.
Контекст достаточен, если покрыты все части:
\begin{equation}
  S(E_t) = \prod_{i=1}^{m} c_t(p_i) = 1.
  \label{eq:sufficiency}
\end{equation}
Определение~\eqref{eq:sufficiency} переносит понятие достаточного контекста, который содержит всю информацию для определённого ответа на вопрос~\cite{joren2024sufficient}, % C142
на уровень отдельных частей вопроса: непокрытые части $U_t = \{p_i : c_t(p_i) = 0\}$ указывают, что именно искать на следующем раунде.
Истинные значения $c_t(p_i)$ системе недоступны: в методе их оценивает языковая модель (раздел~\ref{sec:method}), поэтому далее $\hat{c}_t(p_i)$, $\hat{S}(E_t)$ и $\hat{U}_t$ обозначают оценки этих величин.

\subsection{Критерии качества и стоимости}

Пусть $h(q) = 1$, если ответ на вопрос содержится в корпусе, и $h(q) = 0$ иначе.
Исход оценивается по шкале бенчмарка CRAG, в которой неверный ответ штрафуется, верный поощряется, а пропуск получает ноль~\cite{yang2024crag}: % C247
\begin{equation}
  r(y, q) =
  \begin{cases}
    1, & y \text{ --- верный ответ},\\
    0, & y = \bot,\\
    -1, & y \text{ --- неверный ответ}.
  \end{cases}
  \label{eq:score}
\end{equation}
Среднее значение $r$ на наборе вопросов $Q$ равно $\mathrm{Acc} - \mathrm{Hall}$, где $\mathrm{Acc}$ и $\mathrm{Hall}$ --- доли верных и неверных ответов.
Корректность отказа --- доля вопросов, на которых решение «отвечать или отказаться» совпадает с наличием ответа в корпусе:
\begin{equation}
  \mathrm{AC} = \frac{1}{|Q|} \sum_{q \in Q} \mathbb{1}\bigl[\, \mathbb{1}[y = \bot] = 1 - h(q) \,\bigr].
  \label{eq:ac}
\end{equation}
Величина $\mathrm{AC}$ зависит от доли вопросов без ответа: политика, которая никогда не отказывается, получает $\mathrm{AC}$, равную доле вопросов с ответом.
Поэтому $\mathrm{AC}$ дополняется полнотой отказа на вопросах без ответа $\mathrm{AR}$ и долей избыточных отказов на вопросах с ответом $\mathrm{OA}$:
\begin{equation}
  \mathrm{AR} = \frac{|\{q : y = \bot,\ h(q) = 0\}|}{|\{q : h(q) = 0\}|}, \quad
  \mathrm{OA} = \frac{|\{q : y = \bot,\ h(q) = 1\}|}{|\{q : h(q) = 1\}|}.
  \label{eq:ar-oa}
\end{equation}
Отказом $y = \bot$ считается и отказ по бюджету с обоснованием $\rho$, и ответ «не знаю», который дали генератор или критик (обоснования у него нет).
Ответ «вопрос некорректен» считается ответом и оценивается моделью-судьёй, как любой другой.
Стоимость обработки вопроса описывают число раундов поиска $T(q)$, задержка $L(q)$ без учёта модели-судьи и число токенов $K(q)$.

\subsection{Задача}

Требуется построить политику $\pi$ управления поиском, которая по вопросу $q$ выбирает стратегию и запросы каждого раунда, решает, продолжать ли поиск, и выбирает исход $y$.
Политика должна повышать среднее значение $r$ и корректность отказа $\mathrm{AC}$ относительно базового RAG $\pi_0$, выполняющего один раунд поиска, при ограничении
\begin{equation}
  T_{\pi}(q) \le B \quad \forall q \in Q,
  \label{eq:budget}
\end{equation}
где $B$ --- бюджет итераций.
Относительный рост средней задержки $\bar{L}_{\pi} / \bar{L}_{\pi_0}$ и числа токенов $\bar{K}_{\pi} / \bar{K}_{\pi_0}$ измеряется и входит в результат наравне с качеством.

На решение наложены ограничения:
\begin{itemize}[nosep]
  \item модели не обучаются: поведение компонентов задаётся промптами;
  \item поиск идёт только по закрытому корпусу $\mathcal{D}$, без обращения к внешним источникам;
  \item базовый и агентный конвейеры используют одинаковые корпус, модель эмбеддингов, $k$, число новых фрагментов $k_g$, добавляемых в контекст за раунд, и промпт ответа; за несколько раундов контекст агентного конвейера может вырасти до $N_{\max}$ фрагментов, и это различие входит в цену агентности;
  \item модель-судья, оценивающая ответы, отличается от модели-генератора;
  \item память сессии существует в пределах одного вопроса.
\end{itemize}

\subsection{Схема проверки}

Проверка работоспособности проводится сначала на бенчмарке CRAG в задаче Task 3, где на вопрос приходится до 50 веб-страниц~\cite{yang2024crag}, % C245
затем на собственном наборе вопросов по технической документации и подвыборках TechQA и EManual из RAGBench~\cite{friel2024ragbench}. % C107, C104
Качество ответов оценивается долями верных, неверных и пропущенных ответов и средним $r$, отказы --- величинами $\mathrm{AC}$, $\mathrm{AR}$ и $\mathrm{OA}$.
Качество поиска оценивается метриками Recall@10, Precision@10, MRR@10 и NDCG@10 по первым десяти фрагментам первого запроса; в используемой постановке CRAG разметки релевантных фрагментов нет, поэтому релевантным считается фрагмент, содержащий эталонный ответ, и эти метрики являются прокси-оценками.
Вклад компонентов оценивается абляциями: в каждой из них отключается один компонент --- планировщик, оценка достаточности, критик или память.
Для разностей между системами строятся 95-процентные доверительные интервалы парным бутстрэпом.
Оценки модели-судьи дополняются ручной проверкой спорных случаев.

% ============================================================
\section{Предлагаемый метод}\label{sec:method}

\subsection{Общая схема}

Конвейер реализован как граф состояний из шести узлов: планировщик, поиск, оценка достаточности, генерация, критик и отказ.
Каждый раунд проходит путь «планировщик $\rightarrow$ поиск $\rightarrow$ оценка достаточности».
Если контекст достаточен, управление переходит к генерации и затем к критику.
Если нет и бюджет не исчерпан, начинается следующий раунд с планировщика; если бюджет исчерпан, система переходит к узлу отказа.
Общее состояние графа хранит части вопроса, выполненные запросы, просмотренные фрагменты, накопленный контекст, покрытие частей и номер раунда.
% TODO: рисунок 1 — схема графа (draw.io или TikZ):
% START → planner → retrieve → sufficiency → {generate → critic → END | planner | refuse → END}.

\subsection{Агент-планировщик}

На первом раунде планировщик разбивает вопрос на части~\eqref{eq:parts}; на следующих раундах части не меняются.
Планировщик получает вопрос, время вопроса, список частей, описания того, чего не хватает для непокрытых частей, и уже выполненные запросы.
Для очередного раунда он по списку непокрытых частей выбирает одну из трёх стратегий:
\begin{itemize}[nosep]
  \item \textit{прямой поиск} --- запросом служит вопрос или непокрытая часть без изменений;
  \item \textit{переформулирование} --- короткий запрос из ключевых слов (сущность, атрибут, время) без лишних деталей вопроса;
  \item \textit{декомпозиция} --- отдельный целевой запрос на каждую непокрытую часть.
\end{itemize}
Планировщик формирует от одного до трёх новых запросов только для непокрытых частей.
Запросы, совпадающие с выполненными ранее, отбрасываются.
Если новых запросов не осталось, используется запасной запрос: исходный вопрос, если по нему ещё не искали, иначе вопрос, дополненный описанием недостающих сведений.

\subsection{Поиск и память сессии}

Память сессии $M_t = (Q_t, X_t)$ состоит из множества выполненных запросов $Q_t$ и множества просмотренных фрагментов $X_t$.
На раунде $t$ для каждого запроса $s \in S_t$ выполняется поиск $R(s, k, X_{t-1})$, так что уже просмотренные фрагменты не возвращаются повторно.
Результаты всех запросов раунда объединяются; для фрагмента, найденного несколькими запросами, сохраняется наибольшая оценка сходства.
В контекст добавляются не более $k_g$ новых фрагментов с наибольшими оценками:
\begin{equation}
  E_t = \operatorname{top}_{N_{\max}}\Bigl(E_{t-1} \cup \operatorname{top}_{k_g}\bigl(H_t \setminus E_{t-1}\bigr)\Bigr),
  \label{eq:context}
\end{equation}
где $H_t = \bigcup_{s \in S_t} R(s, k, X_{t-1})$ --- найденные на раунде фрагменты, а $\operatorname{top}_n$ оставляет $n$ фрагментов с наибольшей оценкой сходства.
После раунда память обновляется: $X_t = X_{t-1} \cup H_t$, $Q_t = Q_{t-1} \cup S_t$.

\subsection{Оценка достаточности контекста}

Модуль оценки достаточности получает вопрос, его части и накопленный контекст, в котором у каждого фрагмента есть идентификатор.
Для каждой части $p_i$ модель возвращает оценку покрытия $\hat{c}_t(p_i)$, идентификаторы подтверждающих фрагментов и, если часть не покрыта, краткое описание того, чего не хватает.
Решение принимается только по контексту, без собственных знаний модели.
Часть считается покрытой, только если контекст явно содержит нужные сведения; для частей, зависящих от времени, сведения должны быть верны на момент вопроса.
Часть, основанная на ложной предпосылке, считается покрытой, если контекст показывает, что предпосылка ложна.
Описания недостающего передаются планировщику и определяют запросы следующего раунда.

\subsection{Правило завершения и обоснованный отказ}

После оценки достаточности следующий узел графа выбирается правилом
\begin{equation}
  \operatorname{next}(t) =
  \begin{cases}
    \text{генерация}, & \hat{S}(E_t) = 1,\\
    \text{планировщик}, & \hat{S}(E_t) = 0,\ t < B,\\
    \text{отказ}, & \hat{S}(E_t) = 0,\ t = B.
  \end{cases}
  \label{eq:route}
\end{equation}
Правило~\eqref{eq:route} гарантирует выполнение ограничения~\eqref{eq:budget}.
Простой вопрос, контекст которого признан достаточным после первого раунда, обрабатывается за один раунд поиска; вызовы планировщика, модуля достаточности и критика при этом сохраняются, поэтому и такой вопрос обходится дороже, чем в базовом RAG.
При отказе система возвращает ответ «не знаю» и обоснование $\rho$ --- перечень непокрытых частей $\hat{U}_B$ с описанием недостающих сведений.
В результатах исходы различаются: \textit{answered} (дан ответ), \textit{abstained} (от ответа отказались генератор или критик) и \textit{refused} (бюджет исчерпан до достижения достаточности).

\subsection{Генерация и агент-критик}

Ответ генерируется тем же промптом, что и в базовом RAG: модель отвечает только по приведённым фрагментам, а при нехватке сведений отвечает «не знаю».
Генератору передаётся весь накопленный контекст (до $N_{\max}$ фрагментов); фрагменты, на которые сослался модуль оценки достаточности, ставятся в его начало.
Агент-критик разбивает ответ на атомарные фактические утверждения и для каждого проверяет, подтверждается ли оно фрагментами контекста.
Если все утверждения подтверждены, ответ принимается.
Иначе критик формирует исправленный ответ, в котором остаются только подтверждённые утверждения; если таких не осталось, ответом становится «не знаю».
Ответы «не знаю» и «вопрос некорректен» критик не проверяет.

Работа конвейера сведена в алгоритм~\ref{alg:agentic}.

\begin{algorithm}[t]
  \caption{Агентный цикл обработки вопроса}
  \label{alg:agentic}
  \small
  \begin{algorithmic}[1]
    \Require вопрос $q$, бюджет $B$, параметры $k$, $k_g$, $N_{\max}$
    \Ensure ответ $a$ или отказ $\bot$ с обоснованием $\rho$
    \State $P, U, Q, X, E \gets \varnothing$; $t \gets 0$
    \Loop
      \State $(P, \sigma, S) \gets \textsc{Planner}(q, P, U, Q)$
      \State $H \gets \bigcup_{s \in S} R(s, k, X)$
      \State обновить $E$ по формуле~\eqref{eq:context}
      \State $X \gets X \cup H$; $Q \gets Q \cup S$; $t \gets t + 1$
      \State $(c, U) \gets \textsc{Sufficiency}(q, P, E)$
      \If{$U = \varnothing$} \textbf{выйти из цикла} \EndIf
      \If{$t = B$} \Return $\bot$, $\rho \gets U$ \EndIf
    \EndLoop
    \State $a \gets G(q, E)$
    \State $(v, a') \gets \textsc{Critic}(q, a, E)$
    \If{$v = \text{unsupported}$} $a \gets a'$ \EndIf
    \State \Return $a$
  \end{algorithmic}
\end{algorithm}

Здесь $\sigma$ --- выбранная планировщиком стратегия, $S$ --- запросы раунда, $v$ --- вердикт критика, $a'$ --- исправленный ответ.

\subsection{Реализация}

Конвейер реализован на Python с использованием LangChain и LangGraph: узлы графа соответствуют компонентам метода, а переходы после оценки достаточности задаются правилом~\eqref{eq:route}.
Фрагменты получены разбиением документов на части по 1000 символов с перекрытием 150 символов и хранятся в векторной базе Qdrant; эмбеддинги вычисляются локально моделью BAAI/bge-small-en-v1.5.
Параметры поиска одинаковы для базового и агентного конвейеров: $k = 10$, $k_g = 5$; в агентном конвейере $N_{\max} = 15$ и бюджет $B = 3$.
Базовый RAG выполняет один поиск с $k = 10$ и передаёт генератору пять лучших фрагментов.
Языковые модели вызываются через OpenAI-совместимый API; планировщик, модуль достаточности и критик используют ту же модель, что и генератор, и возвращают структурированный вывод через вызов функций.
Температура равна нулю, начальное значение генератора случайных чисел фиксировано.
Модель-судья относится к другому семейству, чем генератор; конкретные модели фиксируются в плане анализа до прогонов на тестовой выборке.
Для абляций каждый компонент отключается флагом конфигурации: без планировщика запросом служит вопрос, дополненный описанием недостающего; без оценки достаточности все части считаются покрытыми после первого раунда; без критика ответ генератора не проверяется; без памяти просмотренные фрагменты не исключаются из поиска.

% ============================================================
\section*{Декларация об использовании ИИ}

% Скопировано из paper/tex/generated/ai_declaration.tex основного проекта (scripts/build_ai_declaration.py). Руками не править.
При подготовке работы автор использовал Claude Code (Sonnet 5.5 (claude-sonnet-5-5), 31.10.2026) для поиска научной литературы; Claude Code (Opus 5.5 (claude-opus-5-5), 04.11.2026) для поиска научной литературы, помощи в структурировании разделов, подготовки черновых формулировок по заметкам автора и написания программного кода экспериментов; Claude Code (Sonnet 5.5 (claude-sonnet-5-5), агент reader, 04.11.2026) для конспектирования источников, а также Claude Code (Sonnet 5.5 (claude-sonnet-5-5), агент claim-checker, 04.11.2026) для проверки соответствия утверждений источникам. После этого автор проверил и отредактировал текст и несёт полную ответственность за содержание публикации.


% ============================================================
% BEGIN BIBLIOGRAPHY (scripts/bib_to_tex.py)
\begin{thebibliography}{99}

\bibitem{gao2023retrieval} Gao~Y., Xiong~Y., Gao~X. et al. Retrieval-Augmented Generation for Large Language Models: A Survey // arXiv preprint arXiv:2312.10997. — 2023. — DOI: \url{https://doi.org/10.48550/arXiv.2312.10997}.

\bibitem{friel2024ragbench} Friel~R., Belyi~M., Sanyal~A. RAGBench: Explainable Benchmark for Retrieval-Augmented Generation Systems // arXiv preprint arXiv:2407.11005. — 2024. — DOI: \url{https://doi.org/10.48550/arXiv.2407.11005}.

\bibitem{asai2023selfrag} Asai~A., Wu~Z., Wang~Y. et al. Self-RAG: Learning to Retrieve, Generate, and Critique through Self-Reflection // arXiv preprint arXiv:2310.11511. — 2023. — DOI: \url{https://doi.org/10.48550/arXiv.2310.11511}.

\bibitem{joren2024sufficient} Joren~H., Zhang~J., Ferng~C.-S. et al. Sufficient Context: A New Lens on Retrieval Augmented Generation Systems // arXiv preprint arXiv:2411.06037. — 2024. — DOI: \url{https://doi.org/10.48550/arXiv.2411.06037}.

\bibitem{yan2024corrective} Yan~S.-Q., Gu~J.-C., Zhu~Y., Ling~Z.-H. Corrective Retrieval Augmented Generation // arXiv preprint arXiv:2401.15884. — 2024. — DOI: \url{https://doi.org/10.48550/arXiv.2401.15884}.

\bibitem{jeong2024adaptive} Jeong~S., Baek~J., Cho~S. et al. Adaptive-RAG: Learning to Adapt Retrieval-Augmented Large Language Models through Question Complexity // Proceedings of the 2024 Conference of the North American Chapter of the Association for Computational Linguistics: Human Language Technologies (Volume 1: Long Papers). — Association for Computational Linguistics, 2024. — P.~7036–7050. — DOI: \url{https://doi.org/10.18653/v1/2024.naacl-long.389}.

\bibitem{ferrazzi2026agentic} Ferrazzi~P., Cvjetićanin~M., Piraccini~A., Giannuzzi~D. Is Agentic RAG worth it? An experimental comparison of RAG approaches // Proceedings of the 64th Annual Meeting of the Association for Computational Linguistics (Volume 6: Industry Track). — Association for Computational Linguistics, 2026. — P.~56–75. — DOI: \url{https://doi.org/10.18653/v1/2026.acl-industry.5}.

\bibitem{asl2025fairrag} Asl~M.~A., Asgari-Bidhendi~M., Minaei-Bidgoli~B. FAIR-RAG: Faithful Adaptive Iterative Refinement for Retrieval-Augmented Generation // arXiv preprint arXiv:2510.22344. — 2025. — DOI: \url{https://doi.org/10.48550/arXiv.2510.22344}.

\bibitem{li2026s2grag} Li~M., Zou~J., Lv~X. et al. S2G-RAG: Structured Sufficiency and Gap Judging for Iterative Retrieval-Augmented QA // arXiv preprint arXiv:2604.23783. — 2026. — DOI: \url{https://doi.org/10.48550/arXiv.2604.23783}.

\bibitem{yang2025knowing} Yang~D., Zeng~L., Rao~J., Zhang~Y. Knowing You Don’t Know: Learning When to Continue Search in Multi-round RAG through Self-Practicing // Proceedings of the 48th International ACM SIGIR Conference on Research and Development in Information Retrieval. — ACM, 2025. — P.~1305–1315. — DOI: \url{https://doi.org/10.1145/3726302.3730018}.

\bibitem{li2026par2} Li~X., Wang~R., Wang~Y. et al. PAR$^2$-RAG: Planned Active Retrieval and Reasoning for Multi-Hop Question Answering // arXiv preprint arXiv:2603.29085. — 2026. — DOI: \url{https://doi.org/10.48550/arXiv.2603.29085}.

\bibitem{ma2023query} Ma~X., Gong~Y., He~P. et al. Query Rewriting in Retrieval-Augmented Large Language Models // Proceedings of the 2023 Conference on Empirical Methods in Natural Language Processing. — Association for Computational Linguistics, 2023. — P.~5303–5315. — DOI: \url{https://doi.org/10.18653/v1/2023.emnlp-main.322}.

\bibitem{wang2024searching} Wang~X., Wang~Z., Gao~X. et al. Searching for Best Practices in Retrieval-Augmented Generation // Proceedings of the 2024 Conference on Empirical Methods in Natural Language Processing. — Association for Computational Linguistics, 2024. — P.~17716–17736. — DOI: \url{https://doi.org/10.18653/v1/2024.emnlp-main.981}.

\bibitem{singh2025agentic} Singh~A., Ehtesham~A., Kumar~S. et al. Agentic Retrieval-Augmented Generation: A Survey on Agentic RAG // arXiv preprint arXiv:2501.09136. — 2025. — DOI: \url{https://doi.org/10.48550/arXiv.2501.09136}.

\bibitem{jiang2023active} Jiang~Z., Xu~F., Gao~L. et al. Active Retrieval Augmented Generation // Proceedings of the 2023 Conference on Empirical Methods in Natural Language Processing. — Association for Computational Linguistics, 2023. — P.~7969–7992. — DOI: \url{https://doi.org/10.18653/v1/2023.emnlp-main.495}.

\bibitem{es2024ragas} Es~S., James~J., Espinosa Anke~L., Schockaert~S. RAGAs: Automated Evaluation of Retrieval Augmented Generation // Proceedings of the 18th Conference of the European Chapter of the Association for Computational Linguistics: System Demonstrations. — Association for Computational Linguistics, 2024. — P.~150–158. — DOI: \url{https://doi.org/10.18653/v1/2024.eacl-demo.16}.

\bibitem{saadfalcon2024aresautomated} Saad-Falcon~J., Khattab~O., Potts~C., Zaharia~M. ARES: An Automated Evaluation Framework for Retrieval-Augmented Generation Systems // Proceedings of the 2024 Conference of the North American Chapter of the Association for Computational Linguistics: Human Language Technologies (Volume 1: Long Papers). — Association for Computational Linguistics, 2024. — P.~338–354. — DOI: \url{https://doi.org/10.18653/v1/2024.naacl-long.20}.

\bibitem{yang2024crag} Yang~X., Sun~K., Xin~H. et al. CRAG -- Comprehensive RAG Benchmark // arXiv preprint arXiv:2406.04744. — 2024. — DOI: \url{https://doi.org/10.48550/arXiv.2406.04744}.

\end{thebibliography}
% END BIBLIOGRAPHY

\end{document}
